% =========================================================================
% Capítulo 3: Cinética de sistemas de partículas
% Tema 01: Sistemas de partículas
% Problema 04: Sistema de correderas ortogonales acopladas
% =========================================================================
\subsubsection{Sistema de correderas ortogonales acopladas}

Una partícula de masa $m_A$ se desliza sin fricción a lo largo de una barra rígida vertical fija. Esta partícula está unida, a través de un brazo rígido y sin masa de longitud $l$, a una segunda partícula de masa $m_B$, la cual se desliza sin fricción a lo largo de una barra rígida horizontal fija. Suponiendo que $\theta$ es el ángulo formado entre la barra horizontal y el brazo, y que la gravedad actúa hacia abajo, realice lo siguiente:

\begin{enumerate}
    \item Defina los marcos de referencia y sistemas de coordenadas adecuados.
    \item Analice la cinemática de las partículas $A$ y $B$, determinando sus vectores de posición y velocidad en función de $\theta$.
    \item Analice la cinética del sistema para determinar la ecuación diferencial del movimiento en términos del ángulo $\theta$.
\end{enumerate}

\begin{figure}[!ht]
    \centering
    \begin{tikzpicture}[>=Latex, scale=1.5]
        % --- Parámetros ---
        \def\ang{40}      % Ángulo theta
        \def\Lrod{3.5}    % Longitud de la barra l
        \def\RailLen{4.0} % Longitud de los rieles

        % Coordenadas
        \coordinate (O) at (0,0);
        % Posición de B (Horizontal)
        \coordinate (B) at ({\Lrod*cos(\ang)}, 0);
        % Posición de A (Vertical) - Nota: La barra baja, coordenada y negativa
        \coordinate (A) at (0, {-\Lrod*sin(\ang)});

        % --- 1. Rieles y Soportes ---
        % Riel Horizontal
        \draw[line width=4pt] (0, -\RailLen)--(0,0) -- (\RailLen, 0);
        \draw[line width=3pt, gray!40] (0, -\RailLen)--(0,0) -- (\RailLen, 0);
        \draw[] (0, -\RailLen)--(0,0) -- (\RailLen, 0);
        % Soporte pared derecha
        \fill[pattern=north east lines] (\RailLen, -0.2) rectangle (\RailLen+0.2, 0.2);
        \draw[thick] (\RailLen, -0.2) -- (\RailLen, 0.2);

        % Riel Vertical
        %\draw[line width=3pt, gray!40] (0,0) -- (0, -\RailLen);
        %\draw[thick] (0,0) -- (0, -\RailLen);
        % Soporte suelo inferior
        \fill[pattern=north east lines] (-0.2, -\RailLen) rectangle (0.2, -\RailLen-0.2);
        \draw[thick] (-0.2, -\RailLen) -- (0.2, -\RailLen);
        
        % Esquina O
        \node[below right, xshift=2pt] at (O) {$O$};

        % --- 2. Brazo de conexión ---
        \draw[line width=3pt] (A) -- (B) node[midway, above left] {$l$};
        \draw[line width=2pt, gray!40] (A) -- (B);

        % --- 3. Masas (Correderas) ---
        % Masa B (Horizontal) - Dibujada como elipse achatada
        \filldraw[fill=gray!20, draw=black, thick] (B) ellipse (0.2 and 0.2);
        \node[above left, yshift=5pt] at (B) {$m_B$};
        \node[below right, yshift=-5pt, font=\footnotesize] at (B) {$B$};

        % Masa A (Vertical) - Dibujada como elipse achatada verticalmente
        \filldraw[fill=gray!20, draw=black, thick] (A) ellipse (0.2 and 0.2);
        \node[above left, xshift=-5pt] at (A) {$m_A$};
        \node[below right, xshift=5pt, font=\footnotesize] at (A) {$A$};

        % --- 4. Ángulo theta ---
        % Arco desde la horizontal (izquierda de B) hacia la barra
       % \draw[dashed, gray] (O) -- (B); 
        \draw[->, thick] ($(B)+(-0.8,0)$) arc (180:{180+\ang}:0.8);
        \node at ($(B)+(-1.1, -0.3)$) {$\theta$};

        % --- 5. Gravedad ---
        \draw[->, thick] (3, -2.0) -- (3, -3.0) node[midway, right] {$\mathbf{g}$};

    \end{tikzpicture}
    \caption{Sistema de dos partículas deslizantes conectadas por una barra rígida de longitud $l$.}
    \label{fig:correderas_ortogonales}
\end{figure}